\documentclass[12pt]{article}
\usepackage[utf8]{inputenc}
\usepackage[T1]{fontenc}
\usepackage{hyperref}
\title{Yurok Language Resources for K-12 Curriculum Development}
\author{Compiled by Claude Sonnet 5.5 (claude-sonnet-5-5) for Andr\'as Kornai}
\date{}
\begin{document}
\maketitle

\section{Introduction}
Yurok (ISO 639-3 yur, Glottocode yuro1248) is an Algic language of northern California. It is classed as dormant, with no first-language speakers, but it is being revitalized. A 2013 report cited in the sources counts more than 300 basic speakers, 60 intermediate, 37 advanced and 17 conversationally fluent. This report rests on a single search and covers news and programme descriptions only.

\section{Primary Educational Resources}
Reports describe Yurok being taught in schools in northern California. As of 2014 these were four high schools and two elementary schools. The Yurok Tribe's ``Growing Learners and Sustaining Yurok Language'' programme, funded with \$350,000 in fiscal year 2022 by the U.S. Office of Indian Education, covers dual enrollment, digital curriculum and curriculum alignment. It aims to strengthen Yurok courses in four high schools and to double enrolment.
\begin{itemize}
\item CSBA article on the programme: \url{https://blog.csba.org/revitalizing-native-language/}
\item Indianz.com report: \url{https://indianz.com/News/2013/008450.asp}
\item North Coast Journal: \url{https://www.northcoastjournal.com/NewsBlog/archives/2013/02/07/speaking-yurok-in-the-lat}
\end{itemize}
No published curriculum document or textbook was found in the results.

\section{Linguistic and Reference Materials}
None were found in this search.

\section{Religious and Community Resources}
The Yurok Tribe leads the revitalization effort, as described in the articles above. No religious materials were found.

\section{Digital and Audio Resources}
None were found. The programme's digital curriculum is mentioned but no link was located.

\section{Curriculum Development Context}
The language has a funded, school-linked programme run by the tribe. The next step is to contact the Yurok Tribe's education and language department for the actual materials, since none were found online here.
\end{document}